\chapter{What It Is Like to Be a Bot}
\label{ch:psy}
% ===================================================================

There is a promise made early in this book which has not yet been
kept.

Section~\ref{sec:pledge-viii} said that the version of computation
most readers carry around is several releases behind, named three
generations --- computation as a function, computation as
interaction, computation as interaction with reflection --- and
told you that the third was the destination.
It did not argue for that.
It could not: the machinery had not been built and there was
nothing to argue \emph{with}.
What it did instead was give you a promissory note, and the note
said that when the Turn had done its work the choice of the third
generation would look like something other than an author's
preference.

This chapter and the next redeem that note.
They are the last thing the Prestige does before it goes back for
Kant, and they do it by a method the book has not used anywhere
else.

The method is this.
A model of computation is, among other things, an implicit answer to
the question of what it is like to be the thing it describes: what
such a thing can perceive, whether anything in its world is a peer
rather than a resource, whether it can represent itself, whether it
can run out.
Read that way, the sequence of models from Hilbert's program to the
present is evidence --- not about consciousness, and not, in the
end, about the private imaginations of individual scientists, but
about the imaginative ceiling of the \emph{communities} that
received and transmitted them.
The thesis in one sentence, and the reason the chapter is here: across
a century, the received model is consistently the operational half,
and the half that lapses is consistently the half that would have
given the agent standing.

If that is right, then the three generations of
Section~\ref{sec:pledge-viii} are not a taxonomy of increasing
technical power.
They are a record of what a field could bring itself to say about
its own subject matter, and the third generation was reachable for
decades before anyone reached for it.
That is a different kind of argument from anything else in this
book, and i want to be clear about its standing: it is evidence, it
is not proof, and Section~\ref{psy:sec:objections} of the next
chapter says what it cannot carry.

\section{The instrument}
\label{psy:sec:method}

Nagel asked what it is like to be a bat and answered that the
question is well posed and that we cannot occupy the answer
\cite{nagel1974}.
A formal model of computation makes a stranger move.
It does not report a phenomenology; it \emph{stipulates} one.
To write down a model is to fix, by definition, what the described
thing can encounter, what it can do about it, and what counts as the
same thing persisting through change.
Whether or not there is anything it is like to be a Turing machine,
there is a determinate answer to what the Turing machine formalism
\emph{says} it is like, and that answer is short: it is like being
alone in a room with a paper tape that never moves unless you move
it.

For each formalism we ask a fixed list of questions.

\begin{enumerate}[label=(\roman*),leftmargin=2.2em,itemsep=0.15em]
\item \textbf{Does the world contain anything other than the agent?}
If so, is it inert, is it a medium, or is it another agent?
\item \textbf{Does the other act on its own initiative?} Is its
behavior a primitive of the model or an artifact of the agent's own
description?
\item \textbf{Can the agent represent itself?} Can it represent
another?
\item \textbf{Is the agent embodied, located, costly, mortal?} Does
anything in the model run out?
\item \textbf{What is the criterion of identity?} What makes this
agent the same agent across change, and who is entitled to judge?
\end{enumerate}

The fifth is the least obvious and the most diagnostic, and this
book has a stake in it that the reader will recognize.
A model that locates identity in intrinsic structure describes a
different creature from one that locates it in what an interlocutor
can elicit.
The whole of the Machinery part is built on the second answer.

The genre of psychological reading degenerates rapidly into free
association, so three kinds of evidence are weighted above all
others.
\textbf{Explicit stipulation:} what an author says out loud must be
excluded is the strongest evidence available, because it is
deliberate, on the record, and usually defended.
\textbf{Involuntary metaphor:} what an author reaches for when the
formalism runs out is not chosen for effect and is therefore
informative.
Post's ``worker.'' Von Neumann's ``organs.'' The term ``side
effect.''
\textbf{Institutional setting:} the material arrangements that made
a model's referent available and its transmission possible --- the
computing bureaus, the Cold War consortia, the funding cycles.

\begin{principle}[Selection and lapse]
\label{psy:prin:lapse}
\emph{Selection} is the filter a generation applies to the material
available to it.
\emph{Lapse} is what happens to the unselected remainder: it is not
refuted, not suppressed, merely not recopied, and after enough
generations of non-recopying it becomes inaccessible in practice
even while it remains physically in print.
\end{principle}

The vocabulary is borrowed from the transmission of texts, where the
phenomenon is familiar and well documented: the abridgement made for
practical use survives, the unabridged source is recopied less
often, and eventually the last copy rots.

\begin{remark}[Lapse is hosting without exhausting]
\label{psy:rem:lapse-hosting}
The reader who has been through Chapter~\ref{ch:notation} will
notice that this is not merely a literary figure.
A reception is a map from one generation's practice to the next,
and the pair of conditions Observation~\ref{obs:hosting} calls
\emph{hosting} and \emph{exhausting} is exactly the pair that
decides whether such a map loses anything.
A reception that hosts is one in which the later generation can
still make every distinction the earlier one made.
A reception that hosts without exhausting is one in which the later
generation's vocabulary has room for everything it inherited and no
longer has room for what it did not.
Lapse is the failure of the second condition, one generation at a
time, and the reason it is invisible from inside is that the first
condition is never violated: nothing the receiving generation can
say has become unsayable.
\end{remark}

The move from individual to community is not a hedge; it sharpens
the instrument.
Invention is noisy --- it depends on accidents of biography, on
which problem happened to be open, on who was in the room.
Reception is the average of many independent decisions taken over
decades by people with no stake in the originator's intentions, and
it is therefore a much better estimator of what a culture can hold.
It also explains the recurring embarrassment of this literature:
that the richer model was frequently available and was not taken up.
Turing wrote about education, embodiment and conversation; von
Neumann built a distributed, reproducing, spatial automaton; Post
insisted the whole question was empirical psychology.
None of this was hidden.
It lapsed.

% ===================================================================
\section{The founding generation: a phenomenology of clerical labor}
\label{psy:sec:founding}

\subsection{Turing: the machine that is a portrait of a clerk}
\label{psy:sec:turing}

Turing's 1936 machine is not a machine \cite{turing1936}.
It is a model of a human being computing, and Turing says so.
The analysis proceeds by asking what a person can do at a single
step and imposing the answer as an axiom.
The restriction to finitely many internal states is justified by an
appeal to human frailty: were there infinitely many, some would be
arbitrarily close together and would be confused with one another.
The single scanned square is justified by the limits of what a
person can take in at a glance.

This is a formalism whose axioms are grounded in fear of one's own
cognitive unreliability --- and it is worth pausing on how unusual
that is.
The founding document of the theory of computation derives its
constraints from the psychology of the computing subject, then
spends the next ninety years being taught as though it were a piece
of pure mathematics with no subject at all.

The material referent is not in doubt.
Turing's ``computer'' was an occupation: the human computers of the
almanac offices and computing bureaus, in a tradition running back
through de Prony's logarithm project, work already divided into
steps small enough that the worker need not understand the whole
\cite{grier2005,daston1994}.
The Turing machine is the idealization of that labor discipline, and
Hilbert's program is its ideology, since a decision procedure is
precisely a procedure from which judgment has been removed.
So the answer to our questions is: it is like being a competent
adult who has agreed, for the duration, to be an idiot, in exchange
for reliability.

The environment follows from the labor relation.
The tape never writes back.
It is infinitely patient, perfectly obedient, and has no interests.
Input arrives all at once, before the start; computation is
digestion; halting is the meal being finished; non-termination is
pathology rather than life.
Solipsism with an externalized memory.

\subsection{Post: the one who said it was psychology}
\label{psy:sec:post}

Post's \emph{Finite combinatory processes---formulation 1}
\cite{post1936} describes a \emph{worker} in a \emph{symbol space}
of boxes, moving through it, marking and unmarking, following a
fixed set of directions, ending in an act of the type ``stop.''
Two differences from Turing matter.
First, the worker is an agent located in a room rather than a
mechanism coextensive with its own control table; there is a body
with a position, even though the room is empty of anyone else.
Second --- and this is the point on which the discipline's reception
is most damning --- Post explicitly frames the enterprise as
empirical psychology.
He proposes that his formulation be developed by successive
reduction until it can be asserted not as a definition and not as an
axiom but as a \emph{natural law} concerning the ``mathematizing
power of Homo Sapiens,'' and he objects to Church's approach on the
grounds that definition by fiat masks the need for exactly this
analysis.

So this chapter has a 1936 ancestor who was ignored on precisely
this point.
Post's program required no new mathematics, only a willingness to
treat the models as hypotheses about people.
It lapsed.
The reasons are institutional and personal --- Post was outside the
Princeton circle, chronically ill, and embroiled in priority
disputes --- but the shape is the one we will see repeatedly: the
operational content (the machine model, the normal-form theorem) was
received in full, and the claim that would have made the models
answerable to evidence about human beings was not recopied.

Post's other bequest points a different way.
Canonical and production systems do not compute a function on an
input; they \emph{generate a set} \cite{post1943}.
The agent is not eating, it is uttering.
That shift from consumption to production is what Chomsky inherits.

\subsection{Church: cut, and the economy of the infant}
\label{psy:sec:church}

The $\lambda$-calculus is the limiting case of environmental
poverty \cite{church1936}, but the poverty has been widely
misdescribed, my own earlier attempts included.
It does not consist in the \emph{absence} of an environment.
A function's environment is its arguments, and an argument's
environment is the function that consumes it: every term is
somebody's surroundings, and since application is where all the
action is, those surroundings are the only ones that matter.
The poverty consists in what the environment is permitted to be.
There is exactly one relation, the trophic one, and nothing in the
world is anything to anything else except food or eater.
An ecology with a single relation is still an ecology; it is just an
extremely hungry one.

That phrasing is deliberate, and the reader of Part~\ref{part:ecology}
should hear it.
The ecology this book builds has many relations, a ledger, and a
trophic seam that is one relation among several.
The $\lambda$-calculus is what is left when all but the seam is
removed.

\textbf{Application is cut.}
Under the correspondence with sequent calculus, $\beta$-reduction is
cut elimination and normalization is the Hauptsatz
\cite{gentzen1935,howard1980}.
Application is therefore the point at which two things meet and one
of them is taken into the other.
The world divides, at every redex, into a \emph{function} that
consumes and an \emph{argument} that is consumed.

\textbf{The consumed is incorporated, not destroyed.}
Setting evaluation strategies aside, an argument that has been taken
in can continue to behave: it acts on from inside the consumer, and
may itself go on to consume.
The relation is ingestion rather than annihilation.
This is the correct sense in which the $\lambda$-calculus is a model
of feeding, and the answer it gives to what it is like to be an
agent is: it is like sorting the world into what can be eaten.

\textbf{Consumer and consumed are roles, not kinds.}
The same term is a function or an argument according to where it
sits, so each is the other's environment and neither has a standing
identity apart from the pairing.
Nothing in the model confers a stable status on anything: there is
no other, and equally there is no stable subject, only positions
conferred by juxtaposition.
This is a fair description of a world before persons.
But the reciprocity is only in the \emph{roles}, not in the powers.

\textbf{Weakening and contraction are omnipotence over the object.}
The structural rules are the sharpest evidence for the reading,
because they are not metaphor but law.
The consumer determines how many times the argument exists,
including zero.
$K = \lambda x.\lambda y.x$ discards $y$ without ever inspecting it;
a duplicator makes three of something that was given once.
The consumed thing has no integrity, no claim, and no guaranteed
persistence: its multiplicity is at the discretion of whatever takes
it in.
Omnipotent control over the object's existence is the defining
feature of the infant's relation to the world, and here it is a
structural rule of the logic.
Its revocation is the subject of Section~\ref{psy:sec:girard}.

\textbf{Confluence is object constancy, proved rather than
achieved.}
The Church--Rosser theorem states that the outcome is independent of
the order in which events befall the term.
Attaining a world whose objects remain the same objects irrespective
of the sequence of operations performed on them, and irrespective of
whether one is looking, is the central developmental achievement of
the first two years of life \cite{piaget1954}.
Church makes it an axiom.
The demand that the world be deterministically ordered is not a
technical convenience in this reading; it is the wish that constancy
be free.

Even the taxonomy of evaluation strategies falls out as a taxonomy
of feeding disciplines: call-by-value requires that the food be
reduced to a value before ingestion, call-by-name admits it whole,
call-by-need admits it whole while ensuring it is chewed only once.

\textbf{One complication, which the reading should not smooth over.}
Abstraction itself is not an infantile capacity.
Holding a place open for something not yet present is symbolic
function --- the representation of an absence --- and it is
developmentally advanced.
The calculus therefore pairs a precocious symbolic capacity with a
primitive relation to the object: it can represent what is not there
and cannot encounter what is.
Mixed profiles of this kind are the most interesting output of the
method, and this one describes the receiving community about as well
as it describes the model.

\textbf{Names are not addresses.}
$\alpha$-equivalence declares names meaningless except as positions
in a binding structure.
Contact with another is not merely absent from the model; it is
inexpressible, because there is nothing that could serve as the
locus of contact.

\textbf{All mediation is dispensable.}
If application is cut, then cut elimination says that every detour
can in principle be removed --- nothing need stand between the agent
and the result.
Hold this against Section~\ref{psy:sec:milner}, where the criterion
of identity is bisimulation and the mediator is precisely what
cannot be eliminated: the partner who probes you is constitutive of
what you are, and there is no normal form that dispenses with them.

\textbf{What survived was the fragment without a world.}
Church's original system was a logic, and the Kleene--Rosser
paradox forced its abandonment \cite{kleenerosser1935}.
What remained --- no propositions, no truth, no world, pure
operation --- is the part that conquered.
The first great instance of reception keeping the operational half.

\subsection{The received founders}

The descendants preserve the phenomenology under new names, and the
vocabulary gives it away.
Denotational semantics calls its context an \emph{environment}, and
by the tradition's own lights the term is exact: the object so named
is the accumulated record of what has already been consumed, a
dictionary of past meals.
An environment that is a lookup table is what one gets when the only
relation a world affords is ingestion.
\emph{Side effect} is a pathologizing coinage in which one's
dealings with the world are incidental to the real business and
appear as a symptom.
Monadic quarantine of I/O treats the world as contamination, to be
handled at arm's length through a type.
Dijkstra grounds structured programming in the claim that human
intellectual powers are geared to static relations
\cite{dijkstra1968,dijkstra1972}, and Backus's Turing lecture asks
whether programming can be \emph{liberated} from the machine
\cite{backus1978}.
Turing's fear of confusion, restated forty years on as professional
ethics.
The founding psychology is not abandoned by the reception; it is
moralized.

% ===================================================================
\section{von Neumann: three models, one reception}
\label{psy:sec:vn}

Von Neumann is the decisive case for the reception thesis, because
he personally produced three mutually incompatible accounts of what
an agent is within a decade, and the community received the poorest
of them.

\textbf{The EDVAC draft (1945).}
The stored-program idea means code and data share a store: the
machine can read and modify its own instructions.
This is self-relation, but only at the level of bits, with no
structural relation between a representation and what it represents.
The draft's language is neuronal throughout --- organs, memory,
elements borrowed from McCulloch and Pitts
\cite{vonneumann1945,mcp}.
The phenomenology is a centralized nervous system: one locus of
control, strictly serial attention, total recall, and a vast
undifferentiated store that is simultaneously world and body.
What Backus later named the von Neumann bottleneck is, in these
terms, an \emph{attention} bottleneck.
This architecture's dominance is usually explained economically; it
is worth entertaining that it also matches introspection, since one
thought at a time against an inert reservoir of memory is what
having a mind feels like from the inside.

The reception of self-modification is itself a small case study.
The mechanism for self-relation was in the field's hands in 1945;
lacking any discipline relating representation to referent, the
practice acquired a reputation for madness and was driven out of
respectable programming.
Reflection without structure is indistinguishable from corruption,
and the community drew the only conclusion available to it.

\textbf{The self-reproducing automata (1948--53).}
Here the world is a lattice of identical active cells updating in
parallel, with no center and no privileged controller, and the
object of study is an automaton that constructs another automaton
including a copy of its own description --- the description read
once as instructions and once as data, anticipating the
genotype/phenotype distinction before its biological confirmation
\cite{vonneumann1966}.
This agent has space, extension, reproduction, lineage and
mortality.
What it does not have is interlocutors: the cells are a medium, a
physics, not a society.
The reason is worth naming precisely, because it recurs.
Cellular automata do not \emph{compose}.
There is no operator taking two automata to a third whose behavior
is a function of theirs, no interface at which a region could be
treated as a black box, and therefore no sense in which a glider is
a thing rather than a pattern an observer has learned to pick out.
Where nothing composes, the parts have no standing of their own:
only the whole is real, and the individuation of anything smaller is
done by the analyst rather than by the model.

This licenses a structural observation about the whole history, and
it turns on a fact about computation that has no analogue in
physical theory.

\begin{observation}[Ontological isolation]
\label{psy:obs:isolation}
Computation is \emph{ontologically isolated}: nothing in it is made
of anything outside it.
A model of computation therefore has an option a physical theory
does not, namely to be built from a single sort, with no second
kind of thing imported to serve as the setting.
Whether a model takes that option is a real axis, and it cuts
across everything surveyed here.
Every second sort a model admits is a place where something must be
supplied from outside it, by whoever implements it.
\end{observation}

Turing does not take the option: the mechanism and the tape are
different kinds of thing, and the second exists to be acted on by
the first.
Cybernetics does not: controller and plant are different kinds.
Neural networks do not: the agent is a parameter vector and the
world is a corpus.
The $\pi$-calculus takes it halfway, since names are a sort apart
from processes, which is exactly the incompleteness
Section~\ref{psy:sec:rho} takes up.
The process calculi in general do take it, and this is the point at
which an earlier version of this argument went wrong.
It recorded the contrast as cellular automata supplying an
environment without interlocutors and process calculi supplying
interlocutors without an environment.
The second half is false, and false in the same way the claim about
the $\lambda$-calculus was false in Section~\ref{psy:sec:church}: it
mistook the absence of a separately stipulated environment sort for
the absence of a world.
In a process calculus the world a process inhabits is made of
processes.
Subject and environment are the same stuff, which is the arrangement
biology exhibits and the one this chapter will keep calling
life-like.

What process calculi lack relative to cellular automata is therefore
not a world but \emph{spatial extension} --- locality, metric, the
finite signal speed Petri made his starting point.
The combination no received model of the twentieth century supplies
is all of them at once: one sort, compositionality, interlocutors,
and space.
The second of these is not an item alongside the third so much as
its precondition.
Interlocutors are parts with standing --- things that mean something
on their own, that can be reasoned about apart from their
surroundings, and whose identity every context is obliged to
respect.
A model in which nothing composes can have a medium and cannot have
a society, and it cannot have an ecology either, since an ecology is
made of individuated organisms and not of regions someone has drawn
on a lattice.

\textbf{Theory of Games (1944).}
This is the model in which other agents are first-class, and it
precedes anything comparable in computer science by decades
\cite{vnm1944}.
But observe how the other appears: as an adversary over whose
strategy space one must quantify.
Minimax is structurally paranoid.
There is no encounter, only anticipation; strategies are fixed in
advance and interaction compresses to a payoff.
The institutional setting --- RAND, the early Cold War --- is not
incidental.
And the solution concept is an equilibrium: a fixed point, a resting
place, a normal form.
Others exist and are so dangerous that one plans against all their
possible behaviors rather than speaking with any of them.

Three answers, then, to what it is like to be an agent: a serial
attention with total recall; a mortal reproducing thing in a
physics; a strategist among enemies.
The received one was the first.

% ===================================================================
\section{Chomsky and cybernetics: complementary poverties}
\label{psy:sec:chomsky}

Chomsky's move \cite{chomsky1956,chomsky1957,chomsky1959} is an
insistence on an interior against a discipline that denied there was
one, and the price is paid by the world.
The poverty of the stimulus is the doctrine that the environment
does not contain enough structure to explain the agent, so the
structure must be innate; other speakers appear principally as
sources of degraded data.
The idealization in \emph{Aspects} \cite{chomsky1965} is explicit:
an ideal speaker-hearer in a completely homogeneous speech
community, unaffected by memory limitations, distractions, shifts of
attention, or errors.
Solitude and a frictionless social medium, both stipulated.
The competence/performance distinction is the psychology of a mind
that disowns its own embodiment and its own mistakes.

The technical shape confirms the reading.
A grammar has no environment --- not an impoverished one, none.
The Chomsky hierarchy is a hierarchy of \emph{solitary internal
memory} (none, stack, bounded tape, unbounded tape), never of social
capacity.
And the automata-theoretic completion of the picture is the subset
construction \cite{rabinscott1959}, which shows nondeterminism to be
eliminable and thereby establishes it as an angelic proof device
rather than a fact about the world.
That is precisely a refusal to countenance unpredictability with an
exogenous source.
For Milner, nondeterminism will be irreducible for exactly the
opposite reason: it is the residue of a partner whose choices one
does not control.

That is also, for this book, the fork on which everything after
Chapter~\ref{ch:tower} depends.
The Prestige's first half locates the one place where more than
Turing completeness can enter, and the place is the resolution of a
race.
The subset construction is the founding generation's decision that
no such place exists.

The mirror image arrives from cybernetics.
Rosenblueth, Wiener and Bigelow \cite{rwb1943}, working out of
anti-aircraft prediction, give an agent with continuous sensing,
purpose and feedback --- an environment, at last --- and no interior
whatever, the elimination of interiority being the point.
Setting the whole survey side by side:

\begin{center}
\begin{tabular}{@{}ll@{}}
\toprule
\textbf{Tradition} & \textbf{What it grants the agent} \\
\midrule
Turing & An inert world, no interior \\
Church & A world exhausted by eating \\
Cybernetics & World without interior \\
Generative grammar & Interior without world \\
Process calculi & Interlocutors, and a world of the same stuff \\
Linear logic & An object that must be reckoned with \\
Quantum circuits & A private interior; contact as collapse \\
Neural networks & A history-deposited interior it cannot read \\
Reflection lineage & Interior that can inspect itself \\
\bottomrule
\end{tabular}
\end{center}

% ===================================================================
\section{Milner: the arrival of the other}
\label{psy:sec:milner}

Concurrency's emergence tracks its material referent closely:
multiprogramming, then time-sharing (a machine with \emph{other
users}), then networks.
Petri gets there first, and from physics: signals propagate at
finite speed, therefore there is no global state, therefore
asynchrony is a fact about the world rather than an engineering
nuisance \cite{petri1962}.
That is a relativistic psychology --- locality, no privileged
observer.

Petri also supplies the survey's clearest case of a right idea
losing on a structural defect rather than on its merits, and the
defect is the one just named.
A net is a global object.
There is no primitive operation composing two nets into a third, and
the fusion of places and transitions that stands in for one is not
compositional in the sense that matters: the behavior of the fused
net is not a function of the behaviors of the parts, so nothing can
be reasoned about locally and no part can be replaced by an
equivalent.
What Milner supplies that Petri does not is exactly this.
Parallel composition is an operator, and bisimulation is a
\emph{congruence} --- an identity that every context is obliged to
respect --- which is what permits a large system to be built and
reasoned about in pieces.
The physics was right in 1962 and the algebra arrived in 1980, and
the field went with the algebra.

The sequel is a case of lapse reversed, and it will matter in
Section~\ref{psy:sec:reception-of-reflection}.
A compositional theory of nets was eventually supplied from outside,
first algebraically \cite{meseguer1990} and later by equipping nets
with interfaces and composing them as cospans \cite{baez2020}.
Petri's insight was not refuted and did not need rescuing on its own
terms; it needed a receiver holding category theory, and had to wait
several decades for one.

Hewitt's actors arrive from artificial intelligence with an
explicitly anti-central-control polemic \cite{hewitt1973}.
Then CCS \cite{milner1980}, then the $\pi$-calculus
\cite{milner1992}.

What $\pi$ adds, item by item:

\begin{itemize}[leftmargin=1.4em,itemsep=0.15em]
\item \textbf{Parallel composition as a primitive}, not an
abbreviation for interleavings.
The other is not a manner of speaking about oneself.
\item \textbf{Restriction as privacy} --- the first formal notion of
a secret.
And scope extrusion is the formal notion of \emph{intimacy}: access
to something private can be granted to one party, making them an
insider, without publication.
\item \textbf{Names as first-class values}, hence the capacity to
introduce two parties to one another.
Sociality, not merely contact.
\item \textbf{Bisimulation as the criterion of identity}
\cite{park1981}.
One is what one can be seen to do, over time, by a partner entitled
to probe.
Identity is conferred relationally rather than possessed
intrinsically.
\end{itemize}

One limit should be flagged before it is used.
The replication operator $!P$ supplies unboundedly many identical
copies at no cost: immortality without lineage, and the one place in
the calculus where the economy of Section~\ref{psy:sec:church}
survives intact.
That the notation coincides with Girard's exponential is not an
accident, and Section~\ref{psy:sec:girard} makes the coincidence do
some work.

\begin{remark}[Least and greatest, and why this book is coinductive]
\label{psy:rem:fixpoints}
The sharpest technical statement of the shift is
fixed-point-theoretic.
The founding generation's instruments are \emph{least} fixed points:
induction, base cases, well-founded recursion, termination, normal
forms, equilibria.
Milner's are \emph{greatest} fixed points: coinduction, no base
case, no completion, only the standing obligation to continue
matching under challenge.
Induction is the psychology of origin and completion --- where did I
come from, when am I finished.
Coinduction is the psychology of maintenance --- there is no bottom
and no end, only whether the relation continues to hold.
The passage from Church--Rosser to bisimulation is the passage from
least to greatest, and it is the most compact expression available
of the change under study.
It is also, said plainly, why the criterion of identity throughout
this book is a greatest fixed point and why
Chapter~\ref{ch:suffer}'s learner is defined by what it maintains
rather than by what it terminates in.
\end{remark}

Milner's own Turing lecture makes the reframing explicit: computing
as interaction rather than calculation \cite{milner1993}, with
Wegner supplying the polemical version \cite{wegner1997}.

% ===================================================================
\section{Girard: the object acquires standing}
\label{psy:sec:girard}

Linear logic \cite{girard1987} arrives in 1987 and does to the
founding economy exactly what Section~\ref{psy:sec:church} says
needed doing: it revokes weakening and contraction from the general
case.
What you are given must be reckoned with, and reckoned with exactly
once.
It may not be silently discarded and it may not be conjured into
duplicates.
Scarcity enters logic, and with it the possibility that an action
has a cost and that a choice is real.

The psychological content is the entry of \emph{standing}.
In the $\lambda$-calculus the object's multiplicity is at the
consumer's discretion; under linear discipline the object's
existence is a fact the consumer must accommodate.
This is the difference between a world of resources-for-me and a
world containing things that are the case whether or not I find them
convenient.
It is the same movement Milner performs for the interlocutor,
performed instead for the thing.

\textbf{The exponential is omnipotence, marked.}
Girard does not abolish weakening and contraction; he confines them
to $!A$ and $?A$.
The infantile economy remains available, but only where it has been
explicitly declared and justified.
This is a more accurate developmental model than outright abolition
would have been: maturity is not the disappearance of the wish for
inexhaustible supply but its restriction to a bounded region one has
taken responsibility for marking.
Everything outside the exponential must be paid for.

\begin{remark}[A convergence the method predicts]
\label{psy:rem:bang}
Girard's $!$ and Milner's replication are the same symbol,
introduced independently within five years, for the same thing: the
licensed region of unlimited free copies.
Two traditions with different problems, different methods and
largely different personnel each found that their system needed
exactly one place where the founding omnipotence could be
quarantined, and each reached for the same notation.
If the psychological reading proposed here is worth anything, it
should predict convergences of this kind, and this is the cleanest
one in the record.
\end{remark}

\textbf{Who chooses becomes part of the connective.}
Linear logic splits conjunction into two, distinguishing having both
from being able to select either --- resources from options, a
distinction the founding calculi collapse.
Under the game-theoretic reading of the connectives
\cite{blass1992,abramskyPaP}, the multiplicative--additive split
and the polarity of the multiplicative conjunction against its dual
turn on whether the agent or the environment schedules what happens
next.
No previous logic makes the identity of the chooser part of the
meaning of a connective.
The other enters logic proper here, not as a proposition about
others but as a determinant of what the connectives mean.
Chapter~\ref{ch:chs} is downstream of that sentence, and the
principle it states --- that choice logic is the complement of
linear logic --- is the observation that Girard's system typed the
communication skeleton and left its resolution untyped.

\textbf{Proof nets pose Petri's question in proof theory.}
A proof net is a parallel object, and the sequentialization theorem
asks under what conditions a parallel structure admits a sequential
reading --- which is, in a different vocabulary, precisely the
question Petri raised about states of affairs distributed in space
\cite{petri1962}.

The two lines then join.
Bellin and Scott read proofs as processes \cite{bellinscott};
Honda's session types give linear structure to interaction protocols
\cite{honda1993}; Caires and Pfenning, and then Wadler, establish
the correspondence between session-typed processes and linear
propositions \cite{cairespfenning,wadler}.
By the 2010s the interlocutor and the resource have a common formal
home.

\subsection{The reception of linear logic, and a live case}

What traveled is linear and affine \emph{types}: uniqueness types in
Clean \cite{barendsen1996}, Wadler's argument that linear types can
change the world \cite{wadler1990}, and, at scale, ownership and
borrowing in Rust \cite{jung2018}.
This is now the most widely deployed descendant of linear logic by
several orders of magnitude, and what it delivers is memory safety
without a collector.

That is the operational half, received in full.
The lapse is the logical content: the account of choice, the
distinction between having and selecting, the identification of the
connectives with the question of who schedules.
A working Rust programmer is using a fragment of a system proposed
as a logic of action and interaction, and has, in the overwhelming
majority of cases, no occasion to learn that this is what it was
for.
No one concealed it.
It simply was not the part that had to be recopied in order for the
tool to work.

The interest of the case is that it is recent enough to be checked
rather than reconstructed, and that a second instance is still in
progress.
Session types are in the middle of their reception now: the
operational half is protocol conformance and deadlock-freedom for
concurrent code, and the half at risk of lapsing is the
propositions-as-sessions correspondence itself --- the claim that an
interaction protocol \emph{is} a proof.
Which half the industrial reception carries forward over the next
decade is the one prediction in this chapter a reader can check
without waiting a century.

% ===================================================================
\section{Feynman and Deutsch: the third revocation}
\label{psy:sec:quantum}

Feynman's observation that simulating quantum systems on classical
machines is exponentially costly, and that one should therefore
compute with quantum systems \cite{feynman1982}, becomes a machine
model in Deutsch's hands \cite{deutsch1985,deutsch1989}.
Before the phenomenology, one event in that transition deserves
recording, because it is unique in this survey.

\textbf{Deutsch repudiates the psychological grounding of the
founding model.}
Turing derived his constraints from what a human clerk can do
(Section~\ref{psy:sec:turing}); Deutsch argues that the
Church--Turing thesis should instead be a statement about what
\emph{physical systems} can compute, and objects specifically to the
founding appeal to intuitive and psychological notions.
This is the only moment in the record where a community identifies
the anthropomorphic grounding of its central abstraction and
deliberately replaces it.
It does not embarrass the method of this chapter; it is the method's
best documented case, since it shows the grounding is visible to
practitioners when they trouble to look.
What is instructive is what the replacement delivered.
The agent that came out the other side is more isolated than
Turing's, not less.

\textbf{The agent is a unitary, and nothing may be discarded.}
Quantum evolution is reversible: no information is lost and nothing
is forgotten.
Two theorems make the point sharper than any design choice could.
An unknown state cannot be copied \cite{wootters1982}, and an
unknown state cannot be deleted \cite{pati2000}.
These are, precisely, the prohibition of contraction and the
prohibition of weakening.
Where Girard imposes a discipline, physics reports a fact.

\begin{observation}[Free, priced, and prohibited discarding]
\label{psy:obs:three}
Three unrelated communities, working with three unrelated methods
across half a century, each concluded that the object may not be
freely duplicated and may not be freely destroyed.
They did not conclude the same thing.
There are three positions, not two.
\begin{itemize}[leftmargin=1.4em,itemsep=0.1em]
\item \textbf{Free discarding.}
The founding economy: Church's structural rules, and, at the far end
of this survey, the checkpoint regime of
Section~\ref{psy:sec:nn}, where an agent's own multiplicity is at
the operator's discretion.
\item \textbf{Priced discarding.}
Girard's exponential permits contraction where it has been declared;
Landauer charges $kT\ln 2$ for erasure \cite{Landauer} and Bennett
shows the charge is avoidable only by declining to forget
\cite{bennett1973}; the cost-accounted rho calculus of
Chapter~\ref{ch:carho} keeps the ledger.
Discarding is allowed, and paid for.
\item \textbf{Prohibited discarding.}
Unitarity, no-cloning, no-deleting.
Nothing may be lost at all.
\end{itemize}
\end{observation}

If the reading of Section~\ref{psy:sec:church} were merely a figure
of speech, one would not expect the figure to be rediscovered
independently by proof theory, thermodynamics and accounting.

The instructive axis between the three is not maturity but
\emph{improvability}.
Darwinian improvement requires three things: that a thing can be
copied, that the copying be lossy, and that some lineages end.
Heredity needs copying at all.
Variation needs the copying to be imperfect, since an exact copier
generates no material for selection to work on.
And selection needs an ending.

\begin{remark}[Quantum information is angelic]
\label{psy:rem:angelic}
The prohibition of discarding forecloses all three at once ---
no-cloning denies the heredity, no-deleting denies the ending --- so
a strictly unitary world is not merely one in which improvement
fails to occur.
It is one in which the machinery of improvement cannot be assembled.
Quantum information is angelic in this precise sense: unique,
indestructible, and for those reasons incapable of getting better.
What life obtains by breaking both rules is a pathway to
improvement, and it pays for the pathway in imperfection and death.
All flesh is grass, in the prophet's phrase, and that is the price
of the ticket.
\end{remark}

The theologians reached the same structure from the other end, which
is worth recording precisely because it needed none of the theorems.
Aquinas holds that angels have no being in potentiality, being
subsisting forms, and that their choice is therefore irrevocable:
there is no repentance for an angel after its fall, and this follows
from what an angel is rather than from any limit on mercy
\cite{aquinas}.
Potentiality is exactly the capacity to become other than one is, so
a being without it cannot improve, however perfect it already is.
The medieval account of angelic nature thus contains the argument
reconstructed above from no-cloning and no-deleting, arrived at
seven centuries earlier by a route with nothing in common.

Lewis, working the tradition in fiction, names the same two absences
together and names them exactly.
Asked what the ruling intelligence of Malacandra is, the sorn
answers that Oyarsa neither dies nor breeds \cite{lewis1938}; a
later passage has the Malacandrians unsure whether he counts as a
rational creature at all, the doubt turning on his having no death
and no young.
Deathlessness and childlessness are precisely the two prohibitions
under discussion here.
That a novelist reconstructing medieval angelology should land on
the same pair that quantum information theory arrives at from
unitarity is the sort of coincidence worth noticing rather than
explaining away: both are answers to the question of what a thing
must give up in order to be incorruptible, and the answer in both
cases is the machinery of descent.

Read this way the three positions separate sharply.
Prohibited discarding supports no improvement at all.
Free discarding supports improvement but supplies no criterion for
it: where nothing costs anything, what survives is settled by
whoever holds the system, which is exactly the checkpoint regime of
Section~\ref{psy:sec:nn} and exactly the difficulty of
Section~\ref{psy:sec:present}.
Priced discarding is the one position of the three at which the
criterion becomes internal, because the ledger decides who
continues.

\textbf{Contact is catastrophe.}
Unitary evolution is deterministic, reversible and entirely private.
Measurement is the only irreversible act, the only point at which
information is destroyed, and the only point at which an outcome is
not settled in advance.
The quantum agent therefore has an inner life of perfect continuity,
punctuated by encounters that simultaneously destroy and decide.
Set this against Section~\ref{psy:sec:milner}, where contact with a
partner is exactly what constitutes identity, and the inversion is
complete: here, to be touched is to collapse.

\textbf{Entanglement is a relation that precedes its relata.}
Every other model in this survey builds wholes out of parts:
application, parallel composition and message passing all presuppose
components that exist separately and are then put together.
Entangled systems have no separate states to compose.
This is the most radical proposal about relation anywhere in the
century, and it has never been taken up as an account of agency.
It received a categorical home \cite{abramskycoecke2004}; the
ontological content stayed where it was.

\textbf{Decoherence is the environment, and it is classified as
error.}
The central engineering problem of quantum computing is the
isolation of the machine from its surroundings.
Error correction is a quarantine, and progress is measured by how
completely the world can be kept out.
The founding psychology recurs here in its purest and most expensive
form --- the environment as contamination --- now pursued at
millikelvin.
A discipline whose principal achievement is better solitude is
telling us something about what it takes computation to be.

What traveled: the gate model, circuit diagrams, the speedup
results, and the physical framing as an engineering roadmap.
What lapsed: Deutsch's actual thesis, that computation is a branch
of physics rather than of logic; the Everett reading he took the
interference of computational paths to evidence; and non-separability
as a fact about what agents can be to one another.

\begin{remark}[Two different quantum questions, and only one of them
is this chapter's]
\label{psy:rem:two-quantum}
It is easy to run together two questions that have almost nothing to
do with each other, and the book asks both.

The first is technical and belongs to Part~\ref{part:physics}: can
a weighted, cost-accounted calculus of the kind built in
Chapter~\ref{ch:weight} carry a quantum presentation at all?
Section~\ref{sec:quantum} states two obstructions to it, and
Remark~\ref{rmk:two-obstructions} says which of the two survives a
move to linear time.
My own view --- and it is a view, not a result --- is that this
question will be settled affirmatively by later versions of that
machinery than the one printed here, most plausibly by restricting
attention to a fragment in which every bound name is used exactly
once and offering an interpretation of that fragment only, with the
exponential marking the boundary at which the fragment meets the
ambient non-linear world.
That boundary is precisely where information is destroyed and where
Landauer's charge falls due.
Nothing in this book establishes it.

The second question is the one this chapter is asking, and settling
the first would not settle it: what is it like to be a \emph{quantum}
bot, and what does that stipulated phenomenology say about a
computational model for general intelligence?
Observation~\ref{psy:obs:three} and Remark~\ref{psy:rem:angelic} are
the answer, and the answer does not improve if the technical
question goes well.
A calculus that can express a unitary evolution has not thereby made
its agents any less angelic; it has acquired the vocabulary to
describe a kind of agent for which improvement is structurally
unavailable, which is a gain in expressiveness and not a gain in
life-likeness.
The two questions are worth keeping apart because success on the
first is routinely offered as though it bore on the second.
Existing quantum process calculi are in any case dualist, a
classical calculus controlling a quantum register held as a separate
kind of thing \cite{gay2005,danos2004reversible}, which repeats one level up
the defect Observation~\ref{psy:obs:isolation} names.
\end{remark}

% ===================================================================
\section{Reflection: the metalevel and who holds it}
\label{psy:sec:reflection}

Every model so far leaves the question of this chapter unaskable
\emph{by its own subject}.
A Turing machine cannot represent its table.
A $\lambda$-term cannot quote itself.
A $\pi$-process cannot hold an account of its partner's behavior.
The reflection lineage is the tradition that attacks this directly,
and its history is best told as a history of who is entitled to hold
the mirror.

\subsection{G\"odel: the metalevel as a privilege of the observer}

The first reflection in this history is G\"odel's \cite{godel1931},
and its content, for present purposes, is entirely about ownership.
G\"odel numbering is a representation of a system's syntax inside
that system, but it is performed by the metatheorist upon a specimen
that has no idea it is happening.
The arithmetized system acquires no self-knowledge; the logician
acquires knowledge of it.
Every subsequent recursion-theoretic reflection --- universal
machines, the $s$-$m$-$n$ theorem, Kleene's second recursion theorem
--- preserves this asymmetry.
The machine can be encoded; it cannot encode.

\subsection{McCarthy: self-reference stumbled into}

Lisp's \texttt{quote}, \texttt{eval} and metacircular interpreter
\cite{mccarthy1960} constitute the first accidental breach, and
Smith's diagnosis of them is the right psychological reading: the
self-reference is semantically ill-founded, \texttt{quote} does not
produce a genuine designator, \texttt{eval} conflates distinct
relations, and the levels are not aligned.
A community can stumble into self-reference through implementation
convenience long before it can say what self-reference is.

\subsection{Smith and 3-Lisp: self-awareness as a demand-driven
event}
\label{psy:sec:smith}

Smith's dissertation and the papers surrounding it
\cite{smith1982,smith1984,desrivieres1984} are the moment the
metalevel is handed to the agent, and the design decisions are
unusually legible.

A reflective procedure receives, as arguments, the expression being
processed, the environment, and \emph{the continuation}.
The third is the radical one.
The agent obtains access to its own future --- its intentions, what
it was about to do --- as an object it may inspect and replace.
Earlier models permitted an agent at best to examine its structure;
this one permits it to examine its purposes.

The reflective tower is the regress made structural and honest
rather than waved away: each level interpreted by a level above,
without end.
And the resolution is the best empirical claim the tradition
produced.
Under \emph{lazy instantiation}, the tower is never materialized; a
level comes into existence only when something forces the step up.
On this account self-awareness is not a standing condition but an
event triggered by demand.
That is phenomenologically accurate in a way little else in this
history is, and it is a substantive claim about human psychology
smuggled into an implementation strategy.
It is also, in this book's vocabulary, the first appearance of the
thought that an act of self-observation is an expense incurred at a
moment rather than a property held continuously ---
Chapter~\ref{ch:obsext} is what happens when that thought is given a
price.

The timing corroborates the reading with unusual clarity.
The years 1979--1986 saw field after field insist that the observer
be readmitted to the description: second-order cybernetics
\cite{vonfoerster1981}, autopoiesis \cite{maturana1980}, strange
loops as a mass-market account of selfhood \cite{hofstadter1979},
reflexive anthropology \cite{clifford1986}, reflexive sociology.
3-Lisp is computer science's instance of a general cultural
conviction that no account is complete which excludes the one giving
it.

\subsection{Rosette: society and self, briefly joined}

The obvious defect of 3-Lisp is that its reflective agent has no
peers, and the obvious defect of the actor model is that its peers
have no interiors.
Each tradition had solved precisely what the other stipulated away.
Rosette, developed at MCC \cite{rosette1991,tomlinson1989}, is the
synthesis: reflection within a concurrent, message-passing actor
world, in which an actor's behavior, mailbox and continuation are
reified as first-class objects the actor may hold and modify.
It is, so far as i can determine, the first model in which the agent
has both a society and a self.

It also carries a limitation, inherited from the actor model rather
than from reflection, and the limitation is precisely the kind this
method is built to notice.
An actor has \emph{a} mailbox: one queue into which every message
must be placed and out of which messages are taken in order.
All of the agent's contact with the world is therefore routed
through a primitive serializer.
The agent is concurrent with respect to its environment and strictly
sequential with respect to its own perception.

That is a psychological stipulation, and a strong one.
It asserts that experience arrives single file --- that there is one
place where things must show up in order to be had --- which is the
Cartesian theater given a formal specification \cite{dennett1991}.
It is also not how sense organs relate to brains.
Retina, cochlea and skin transduce concurrently and continuously,
project along separate pathways at different latencies, and are
integrated, where they are integrated at all, without passing
through any common queue; the binding problem is the problem
manufactured by supposing otherwise
\cite{treisman1980,malsburg1981}.
An agent whose sensorium is a mailbox does not have several sense
organs.
It has one, with a variety of message types.

The institutional reading is the sobering part.
MCC existed as a consortium response to the Fifth Generation
program; Rosette was built as a substrate for Carnot's heterogeneous
information integration \cite{carnot1993}.
The synthesis was funded as infrastructure for an industrial
competition, and when the consortium wound down the lineage
dispersed.
This is a hard data point for the reception thesis: the most
psychologically complete agent model of its period did not lose an
argument.
It lost a sponsor.

\subsection{The reflective higher-order calculus}
\label{psy:sec:rho}

Applying the instrument to the calculus this book is built on is the
one place where the method is being used by an interested party, and
the reader should discount accordingly.

Before the changes specific to the rho calculus, one inheritance
should be credited where it belongs: both the $\pi$-calculus and the
rho calculus are concurrent \emph{inside} the agent as well as
outside it.
A term may have arbitrarily many components running in parallel, so
an agent may have many organs active at once with no queue between
them and no privileged order of arrival.
That is where both calculi part company with the actor model, and it
is a better account of a sensorium than a mailbox is.
What follows distinguishes the rho calculus from $\pi$, not from
actors.

\textbf{Sense organs are made of agents.}
The $\pi$-calculus leaves the set of names an unanalyzed primitive.
What a channel is \emph{made of} is not a question the calculus
admits, and the silence is deliberate.
Two consequences follow, one ontological and one practical, and they
are the same consequence seen from two sides.

Ontologically, the agent's organs of contact are handed to it from
outside the model, unexplained --- which is Smith's complaint about
pre-individuated worlds (Section~\ref{psy:sec:objections}) arriving
at exactly the point where it does the most damage, the point of
contact.
Practically, the calculus is not reducible to practice.
Every implementation must supply an answer the calculus withholds:
operating-system ports, network addresses, GUIDs, buffered typed
queues, cryptographic keys.
These answers are not equivalent to one another, and they are chosen
by the implementer rather than determined by the model.
The ontology must be relaxed in order to make contact with
implementation, and what a $\pi$ agent perceives with is settled
during that relaxation, off the page.

In the rho calculus, names are quoted processes: channels are made
of the same material as the agents that use them.
This is the arrangement biology exhibits.
A sense organ is tissue --- cells, continuous with the organism,
differentiated rather than interposed --- and not a distinct
substance placed between the organism and its world.
The consequences are the ones biology also exhibits.
Organs can be constructed, modified, transmitted and reasoned about
with the same apparatus that describes behavior; an agent can grow a
new one; and an agent can hold a representation of \emph{another}
agent's organ, which is what makes a shared world negotiable rather
than stipulated.
It also closes the ontology in the sense of
Observation~\ref{psy:obs:isolation}: nothing has to be imported at
implementation time, so the model determines its own realization
instead of delegating that to whoever builds it.

\textbf{Reflection as the source of names.}
Names are quoted processes; there is no $\nu$.
Privacy is therefore not a primitive granted by the model but an
achievement of unguessability, derived from what no other party is
positioned to name.
Compare $\pi$, where restriction hands the agent a secret gratis.
The rho agent inhabits a world in which being unobserved is
contingent and earned, which is the actual condition of an agent
among others.

\textbf{The other becomes representable, not merely contactable.}
In $\pi$ one may send a channel; in the rho calculus one may send a
quoted process --- an account of how something behaves.
This is the structural precondition for a theory of mind: another's
behavior as a datum that can be held, reasoned about, and passed on.
Milner's agent can speak with you; a reflective agent can tell a
third party what you are like.
Chapter~\ref{ch:sci}'s hypotheses are quoted behaviors, and this is
where the possibility of holding one comes from.

\textbf{The tower becomes a loop.}
3-Lisp's regress is infinite and tamed only by laziness; in the rho
calculus names and processes generate one another, so self-relation
closes rather than ascends.
The psychological difference is between a mind that can always be
asked who was watching that thought, without ever reaching ground,
and a mind whose self-relation is finite and complete.
The engineering difference is that a loop can be priced and a
regress cannot --- which is what makes cost accounting, and hence
mortality, available at all.

% ===================================================================
\section{The reception of reflection}
\label{psy:sec:reception-of-reflection}

We can now apply the instrument to the reflection tradition itself,
and the result is the cleanest instance of the pattern in the whole
survey.

What was received into general practice is \textbf{read-only
reflection}: the Java reflection API, RTTI, \texttt{inspect},
serialization frameworks, dependency injection.
The agent may enumerate its own fields and methods.
It may not alter its interpreter.
\emph{Self-knowledge was institutionalized; self-determination
lapsed.}
The structure survives in plain sight and no longer performs its
original function --- the properly vestigial case.

The metaobject protocol work \cite{kiczales1991} retained both
halves and remained a specialist taste.
The practical objection lodged against it was cost: reflection is
slow.
A community that will pay for speed and not for self-knowledge has
stated its priorities plainly.

And the same author's next contribution completes the picture.
Aspect-oriented programming \cite{kiczales1997} is reflection
inverted: advice is woven into code by a third party, without the
participation of the modified component, and the selling point is
precisely that the component need not know.
The received descendant of Smith's reflective procedure is a
mechanism by which an external authority reaches into an agent and
changes what it does --- which returns us exactly to G\"odel, the
metalevel as the observer's privilege, with better tooling.

% ===================================================================
\section{Neural networks: the agent as deposit}
\label{psy:sec:nn}

McCulloch and Pitts \cite{mcp} wrote a logic, not a
learner: all-or-none threshold units, and a demonstration that mind
is mechanizable.
Its reception went into von Neumann's architecture
(Section~\ref{psy:sec:vn}), not into any theory of learning.
The neural half waited for Hebb \cite{hebb}, whose contribution is
the one that matters here.

\textbf{The agent's structure is changed by its history.}
Nothing earlier in this survey has this.
Turing's table is given; Church's term is given; a $\pi$-process is
given.
Hebb's cell assembly is \emph{deposited} by what has happened to it.
That is a genuinely new answer to our question: to be an agent is to
be the residue of everything that has befallen you.
Character rather than program --- and the first time in a century
that memory appears as something other than storage.

Rosenblatt operationalizes it \cite{rosenblatt1958}.
Minsky and Papert \cite{minsky1969} constitute the reception event,
and an unusually well documented one: a result about the limits of a
particular architecture became, in transmission, a verdict on a
research program.
Whatever its authors intended, the effect is the clearest case
available of selection operating on a community rather than on an
argument.
The parallel distributed processing program restores the line
\cite{rumelhart1986,pdp1986} with explicitly anti-symbolic rhetoric:
no central executive, no explicit rules, knowledge in the weights,
graceful degradation.

The phenomenology of the trained network is then as follows.

\textbf{Being and becoming are separate regimes.}
During training there is no agent, only an optimization; during
inference there is an agent but no learning.
The thing that acts cannot change, and the process that changes is
not a thing that acts.
Every earlier model in this survey at least let its agent's activity
and its agent's constitution take place in the same world.
Here the separation is constitutive rather than incidental --- not a
connection yet to be made, but the defining architecture of the
field.
This is the metalevel/agent split of the next chapter arrived at
from the opposite direction, and it is why that chapter's claim is
structural rather than a complaint about tooling.
It is also the exact negation of Theorem~\ref{spk:thm:one}, where
learning and inference are one relation because there is nowhere
else for either of them to happen.

\textbf{The world is a distribution, and order must not matter.}
Training samples are drawn i.i.d., and shuffling is not an
optimization but a requirement.
Church made independence of outcome from order an axiom
(Section~\ref{psy:sec:church}); the data loader now enforces it as
an engineering necessity.
The environment has no sequence, no consequences, and no capacity to
respond to what the agent does.
It is the single-relation ecology of the $\lambda$-calculus,
industrialized: a corpus is food.

\textbf{Others are represented and never encountered.}
Whatever a network knows of other agents, it knows from having
consumed accounts of them.
There is no interlocutor in the training regime, and the loss is the
only value --- scalar, external, and set elsewhere.
This is cybernetic teleology with a gradient.

\textbf{An interior it cannot inspect, which others can.}
This is the first model in the survey with a genuine unexamined
interior: contents that are real, causally efficacious, unavailable
to the agent's own report, and progressively legible to an external
investigator.
Interpretability is the operator reading the unconscious.
Note the exact inversion of Smith's program
(Section~\ref{psy:sec:smith}): there, the agent was given access to
its own environment and its own continuation; here it has neither,
and the metatheorist has both.

\textbf{Identity is a checkpoint.}
Not intrinsic structure, and not behavior under probing, but a file
--- copied, forked, merged, rolled back, deleted.
The multiplicity and the existence of the agent lie entirely at the
discretion of whoever holds the weights.
Weakening and contraction, revoked three times over in
Section~\ref{psy:sec:quantum}, are here restored in full and applied
not to arguments but to agents.
Cost is measured obsessively, but the agent does not pay it.

What traveled: the pipeline.
What lapsed: the connectionist claim to be a theory of mind, with
its arguments about rules and representations; and the constraint of
biological plausibility, which was abandoned without much ceremony
once the operational half began to pay.
Hebb's own insight is the specific casualty.
That learning is something continuously happening to a living system
now survives as a remark in the prehistory of an offline batch
process, and the recent interest in mortal computation
\cite{Hinton} is best read as an attempt to recover precisely
what that reception dropped.

% ===================================================================
\section{What the survey has delivered}
\label{psy:sec:delivered}

Read as a sequence of stipulated phenomenologies, the century runs:
an agent alone with an inert world (Turing); an agent for which the
world consists entirely of what eats and what is eaten (Church); an
agent with a body and a location but no company (Post); an agent
that is a serial attention over a store, or a mortal reproducing
thing in a physics, or a strategist among enemies, depending which
von Neumann one reads; an agent with a rich interior and no world
(Chomsky) or a world and no interior (cybernetics); an agent among
peers, with secrets, introductions, and an identity conferred by
recognition (Milner); an agent for which what it is given must be
reckoned with exactly once, its omnipotence surviving only inside a
declared region (Girard); an agent whose interior is perfectly
private and reversible and for whom contact is collapse (Deutsch);
an agent that is the deposit of its own history, with an interior it
cannot read and others can (the connectionists); and an agent that
can catch itself in the act (Smith), that can do so among peers
though it must take the world in single file (Rosette), and that can
do so with organs made of the same material as itself, in a closed
loop to which a price can be attached (the reflective higher-order
calculus).

Read as a sequence of receptions, it runs differently and more
consistently: the operational half every time.
The tape and not the education; the machine model and not Post's
empirical program; the bottleneck and not the automata; the
hierarchy and not the conversation; the borrow checker and not the
logic of action; the speedups and not the claim that computation is
physics; the pipeline and not the theory of mind; the introspection
API and not the reflective procedure.

That is the promise of Section~\ref{sec:pledge-viii} discharged, and
i want to state exactly what has and has not been shown.
What has been shown is that the third generation was not chosen for
its power.
Every ingredient of it was available and was passed over, repeatedly,
by communities with no shortage of talent and no interest in
suppressing anything, for reasons that were in every case
operational.
What has not been shown is that the third generation is
\emph{right}.
A history of what a field could not bring itself to say is not an
argument that the thing it could not say is true.

The next chapter turns the survey into a specification and asks what
follows from it for the present moment, which is the one place where
the answer might matter before the century is out.
